% =============================================================================
%  STUDY COMPANION — Foundations of Conversational and Agentic AI
%  ZG521 Conversational AI, Session 1
%  Built with make-lecture-kit v3.1.0
% =============================================================================

\documentclass[a4paper,11pt]{article}

\usepackage[T1]{fontenc}
\usepackage[utf8]{inputenc}
\IfFileExists{mathpazo.sty}{\usepackage[sc]{mathpazo}\linespread{1.06}}{\usepackage{lmodern}}
\IfFileExists{helvet.sty}{\usepackage[scaled=0.9]{helvet}}{}
\renewcommand{\ttdefault}{lmtt}
\usepackage{microtype}

\usepackage[a4paper,top=26mm,bottom=26mm,left=27mm,right=27mm,
            headheight=14pt,headsep=12mm,footskip=12mm]{geometry}

\usepackage{amsmath,amssymb}

\usepackage{graphicx}
\usepackage{booktabs}
\usepackage{array}
\usepackage[table,svgnames]{xcolor}
\usepackage{caption}
\usepackage{float}
\usepackage[section]{placeins}
\usepackage{needspace}

\usepackage[most]{tcolorbox}
\tcbuselibrary{skins,breakable}

\usepackage{tikz}
\usepackage{pifont}

\usepackage{enumitem}
\setlist{leftmargin=1.5em,topsep=4pt,itemsep=3pt}

\newlist{steps}{enumerate}{1}
\setlist[steps]{label={\sffamily\small\bfseries\color{accent}Step \arabic*},
  leftmargin=4.2em, labelwidth=3.6em, labelsep=0.6em, labelindent=0pt,
  align=left, topsep=6pt, itemsep=5pt}
\usepackage{parskip}
\setlength{\parskip}{0.55em plus 0.1em}

\usepackage{fancyhdr}
\usepackage{titlesec}

\usepackage[hidelinks]{hyperref}
\usepackage{xurl}
\hypersetup{breaklinks=true}

\definecolor{accent}{HTML}{1F4E79}
\definecolor{inktext}{HTML}{1F2328}
\definecolor{mutedtext}{HTML}{6B7280}
\definecolor{hairline}{HTML}{D0D5DD}
\definecolor{wash}{HTML}{F5F6F8}
\colorlet{navy}{accent}      \colorlet{navyrule}{hairline}
\colorlet{intuFrame}{accent} \colorlet{everyFrame}{accent} \colorlet{workFrame}{accent}
\colorlet{watchFrame}{inktext} \colorlet{keyFrame}{accent} \colorlet{waitFrame}{accent}

\color{inktext}
\hypersetup{linkcolor=accent,urlcolor=accent,citecolor=accent}

\newcommand{\SeriesName}{ZG521 Conversational AI}
\newcommand{\TopicTitle}{Foundations of Conversational \& Agentic AI}

\pagestyle{fancy}
\fancyhf{}
\fancyhead[L]{\sffamily\scriptsize\color{mutedtext}\MakeUppercase{\SeriesName}}
\fancyhead[R]{\sffamily\scriptsize\color{mutedtext}\TopicTitle}
\fancyfoot[C]{\sffamily\footnotesize\color{mutedtext}\thepage}
\renewcommand{\headrulewidth}{0pt}
\renewcommand{\footrulewidth}{0pt}
\fancypagestyle{plain}{\fancyhf{}\fancyfoot[C]{\sffamily\footnotesize\color{mutedtext}\thepage}}

\titleformat{\section}
  {\sffamily\Large\bfseries\color{inktext}}
  {\color{accent}\thesection}{0.7em}{}
\titleformat{\subsection}
  {\sffamily\large\bfseries\color{inktext}}
  {\color{accent}\thesubsection}{0.6em}{}
\titlespacing*{\section}{0pt}{26pt}{6pt}
\titlespacing*{\subsection}{0pt}{14pt}{4pt}
\newcommand{\sectionbreak}{\needspace{14\baselineskip}}

\newcommand{\secsub}[1]{%
  \noindent{\large\itshape\color{mutedtext}#1}\par
  \vspace{2pt}{\color{hairline}\rule{\linewidth}{0.4pt}}\par\medskip}

\newcommand{\flipanswer}[1]{\par\noindent\hfill\rotatebox[origin=c]{180}{\parbox{0.9\linewidth}{\small\color{mutedtext}#1}}\par\medskip}

\tcbset{
  quiet/.style 2 args={
    enhanced, breakable, sharp corners,
    colback=#2, colframe=#2, boxrule=0pt,
    borderline west={1.6pt}{0pt}{#1},
    left=12pt, right=10pt, top=7pt, bottom=7pt,
    before skip=12pt, after skip=12pt,
    fontupper=\normalsize, coltitle=#1,
    fonttitle=\sffamily\scriptsize\bfseries,
    attach title to upper={\par\smallskip},
  }
}
\newcommand{\calloutlabel}[1]{\MakeUppercase{#1}}

\newtcolorbox{intuition}{quiet={accent}{white},
  title={\calloutlabel{The intuition}}}
\newtcolorbox{everyday}{quiet={accent}{white},
  title={\calloutlabel{Everyday picture}}}
\newtcolorbox{worked}[1]{quiet={accent}{wash},
  title={\calloutlabel{Worked example}\ \textperiodcentered\ {\normalfont\itshape\small\color{inktext}#1}}}
\newtcolorbox{watchout}{quiet={inktext}{white},
  title={\calloutlabel{Watch out}}}
\newtcolorbox{waitwhat}{quiet={accent}{white},
  title={\calloutlabel{Wait --- really?}}}
\newtcolorbox{keytake}{quiet={accent}{white},
  borderline north={0.4pt}{0pt}{hairline},
  borderline south={0.4pt}{0pt}{hairline},
  title={\upshape\calloutlabel{Key takeaway}}, fontupper=\normalsize\itshape}

\newcommand{\titlebanner}[4]{%
  \begingroup\raggedright
  \vspace*{4mm}
  {\sffamily\footnotesize\bfseries\color{accent}\MakeUppercase{#1}}\par
  \vspace{7pt}
  {\sffamily\bfseries\fontsize{30}{35}\selectfont\color{inktext}#2\par}
  \vspace{7pt}
  {\large\color{mutedtext}#3\par}
  \vspace{12pt}
  {\color{hairline}\rule{\linewidth}{0.5pt}}\par
  \vspace{3pt}
  {\sffamily\scriptsize\color{mutedtext}\MakeUppercase{#4}}\par
  \vspace{10pt}
  \endgroup
}

\captionsetup[figure]{font={small,color=mutedtext},labelfont={bf,sf,color=inktext},
  labelsep=period,justification=raggedright,singlelinecheck=false,skip=6pt}
\setlength{\textfloatsep}{16pt plus 4pt minus 4pt}
\setlength{\intextsep}{14pt plus 4pt minus 4pt}
\renewcommand{\topfraction}{0.8}\renewcommand{\bottomfraction}{0.6}
\renewcommand{\textfraction}{0.15}\renewcommand{\floatpagefraction}{0.75}
\newcommand{\housefig}[3][0.86]{%
  \begin{figure}[!htbp]
    \centering
    \includegraphics[width=#1\linewidth]{#2}%
    \caption{#3}
  \end{figure}%
}

\renewcommand{\arraystretch}{1.18}
\newcommand{\vocab}[1]{\textbf{#1}}

\begin{document}

\thispagestyle{fancy}

\titlebanner
  {\SeriesName}
  {Foundations of Conversational \& Agentic AI}
  {From Eliza to autonomous agents: what actually changed, and why it matters}
  {Built from: the Session 1 handout, syllabus and recording · Every example worked out in full}

\noindent\textbf{How to use this companion.}
This companion explains Session 1 from scratch, in plain English. Every idea
starts with an everyday picture. Then the mechanism is built up one step at a
time. Every example is worked out in full. Small labels mark the boxes.
\textsf{\small\textbf{THE INTUITION}} states the core idea in one breath.
\textsf{\small\textbf{EVERYDAY PICTURE}} ties it to real life before any
mechanism. \textsf{\small\textbf{WORKED EXAMPLE}}, on a grey background,
walks a real task through with no skipped steps. \textsf{\small\textbf{WATCH
OUT}} flags the common mistake. \textsf{\small\textbf{WAIT --- REALLY?}}
holds a result that sounds wrong but is true. \textsf{\small\textbf{KEY
TAKEAWAY}} is the one sentence to remember. Read it in order; each section
leans on the one before.

\medskip

% =============================================================================
\section{Try it first: what would you need to build?}
\secsub{Before any definitions, trace one real request and find out, by
feel, what a simple chatbot cannot do.}

Here is a request a user types into a travel-booking app:

\begin{quote}
\itshape
``Book me a window seat on tomorrow's 7am flight to Chennai, under
Rs.~6{,}000, and text me the confirmation.''
\end{quote}

Before reading on, picture a 1990s-style chatbot trying to answer this. It can
match patterns in your sentence and print a canned reply. Ask yourself three
questions. Can it \emph{check today's flight prices}? Can it \emph{remember}
your Rs.~6{,}000 condition while it checks? Can it \emph{act} --- actually book
the seat and send a text? The answer to all three is no. That gap is exactly
what this session is about.

A system that can only generate a reply is missing three things. First, a way
to look things up or act in the world --- we call this a \vocab{tool}.
Second, a way to hold a condition in mind while it works --- \vocab{memory}.
Third, a way to decide, on its own, what to do next --- \vocab{autonomy}.
Everything below names and builds these three pieces.

% =============================================================================
\section{Four eras, one evolving idea}
\secsub{Four different attempts answer the same question. Each one fixes the
last attempt's biggest weakness.}

Every system in this history answers one question: given what the user just
said, what should the system do next? The four eras differ only in \emph{how}
they answer it.

\begin{intuition}
Each era kept the parts of the previous one that worked. It replaced only the
part that broke. Rules broke on anything unscripted, so retrieval replaced
them. Retrieval broke on anything novel, so a generative LLM replaced it. A
bare LLM broke the moment real-world action was needed, so an agent loop
wrapped around it.
\end{intuition}

\begin{everyday}
Picture four generations of a customer-support desk. Generation 1 is a new
hire reading word-for-word from a script --- fine, until a customer says
something off-script. Generation 2 is a better hire who searches past tickets
for the closest match and copies that reply. Generation 3 is a sharp hire who
has read a huge library. They write a fresh, fluent answer to almost
anything --- but have no phone, no login, and no memory of yesterday's call.
Generation 4 is that same hire, now handed a phone, a login, and a notebook.
They can look things up, take real action, and remember what you said five
minutes ago.
\end{everyday}

\begin{table}[htbp]
\centering\small
\begin{tabular}{@{}p{2.6cm}p{4.0cm}p{4.3cm}@{}}
\toprule
\textbf{Era} & \textbf{How it decides a reply} & \textbf{Where it breaks} \\
\midrule
Eliza, Alice (rule-based) & Hand-written pattern $\to$ fixed reply template &
  Anything outside the written rules \\
Retrieval / ML chatbots & Pick the closest-matching reply from a trained set &
  Anything genuinely new, never seen before \\
LLM assistants & Generate fresh text, one token at a time, from what it
  ``read'' during training & No memory, no tools --- it can only talk, never
  check or act \\
Autonomous agents & Loop: reason, call a tool, read the result, repeat, with
  memory across turns & Reliability and cost --- now the hard problems, not
  fluency (Sections 3--4) \\
\bottomrule
\end{tabular}
\caption*{\footnotesize The four eras, compressed.}
\end{table}

\housefig{figures/fig_evolution.png}{Four eras of the same idea: each one
removes the previous generation's biggest limitation.}

\begin{worked}{tracing one request, ``find me a good Italian place nearby'', through all four eras}
\begin{steps}
  \item \textbf{Era 1 --- rules.} The system matches the word ``Italian'' to a
    pre-written rule: \texttt{IF message contains ``Italian'' THEN reply
    ``Try Giorgio's on Main Street.''} It always gives the same answer. It
    does not matter if Giorgio's is open, nearby, or still in business.
  \item \textbf{Era 2 --- retrieval.} The system has thousands of past
    support replies. It finds the one most similar to your message and
    returns it verbatim. Coverage is better, but it is frozen at training
    time --- it cannot check what's open \emph{right now}.
  \item \textbf{Era 3 --- LLM assistant.} The system reads your message and
    \emph{generates} a fresh, fluent reply: ``Here are a few well-reviewed
    Italian restaurants in your area...'' It sounds great. But it is
    reciting from memory, not checking a live map. It may invent a
    restaurant that closed last year --- a \vocab{hallucination}: a fluent,
    confident, but unverified claim.
  \item \textbf{Era 4 --- agent.} The system reasons: ``I need current,
    nearby results, so I should call a maps search tool, not guess.'' It
    calls that tool with your location. It reads the real results back, and
    only then writes its reply --- grounded in a live lookup, not a guess.
\end{steps}
\end{worked}

\begin{watchout}
Newer does not mean ``always better, no trade-offs.'' An agent calling tools
is slower and costs more than a rule or a cached reply. Every tool call, and
every reasoning step, has a real cost. A spam filter or an FAQ bot often
still uses Era 1 or 2 technology, on purpose. The task is simple, and an
agent loop would be wasted effort. Picking the right era for the task is
itself a design skill you will use all semester.
\end{watchout}

\begin{keytake}
Four eras answer the same question --- what should the system do next? ---
with increasingly capable, increasingly expensive machinery: fixed rules,
then retrieved replies, then freshly generated text, then a
generate-act-observe loop with memory. Each later era fixes the previous
one's sharpest limitation.
\end{keytake}

\noindent\textbf{Where it's used.}
This same four-era ladder shows up across the industry, not just in chat.
Search engines moved from keyword rules, to ranking models, to
retrieval-augmented generation. Customer support moved from IVR menus, to ML
ticket routing, to LLM copilots. Today it can reach fully autonomous agents
that issue refunds on their own. The pattern --- rules, then learned matching, then
generation, then action --- repeats wherever software responds to open-ended
human input.

% =============================================================================
\section{What actually makes something an agent}
\secsub{Five pieces turn a text generator into an agent. A system needs all
five, not just one.}

A language model by itself is extremely good at one thing. Given some text,
it predicts good text to come next. That is \vocab{reasoning} in its rawest
form --- working out what should logically come next. On its own, reasoning
is not enough. A brilliant intern who has never left the office cannot book
your flight either. Being right about what \emph{should} happen is not the
same as being able to \emph{make} it happen.

\begin{intuition}
An agent is not one new invention. It is five existing ideas wired together
around a language model. Four of them are reasoning, planning, tools, and
memory. The fifth is an autonomy setting. That setting decides how much the
agent may do without asking you first.
\end{intuition}

\begin{everyday}
A guide dog is a working definition of ``agent'' with fur. It
\emph{reasons} about the scene ahead (is that a step down, a car, an open
door?). It \emph{plans} a safe route to the destination. It uses its
\emph{tools} --- its body, its nose, its training --- to act in the world. It
\emph{remembers} its handler's routines and the route home. And it has real
\emph{autonomy}. It can refuse the ``forward'' command if a car is coming,
acting against the instruction, because its training says to. Take away any
one of the five and it stops being a guide dog. No planning, and it wanders.
No memory, and it forgets the route every ten metres. No autonomy, and it
walks you into traffic because it was told to.
\end{everyday}

\begin{table}[htbp]
\centering\small
\begin{tabular}{@{}p{2.0cm}p{4.6cm}p{4.3cm}@{}}
\toprule
\textbf{Piece} & \textbf{What it does} & \textbf{In the flight-booking example} \\
\midrule
Reasoning & Decides what needs to happen next, step by step &
  ``I need today's price before I can compare it to Rs.~6{,}000.'' \\
Planning & Breaks one goal into an ordered sequence of actions &
  1. Check price. 2. Compare to budget. 3. Book if under budget. 4. Send text. \\
Tools & The agent's ``hands'': APIs, search, code, databases &
  A flight-price API, a booking API, an SMS API \\
Memory & Carries state across steps and across turns &
  Remembering ``window seat, under Rs.~6{,}000'' while three tool calls run \\
Autonomy & Decides when to retry, adapt, or stop, inside set limits &
  Trying a different flight if the first is over budget, up to a retry limit \\
\bottomrule
\end{tabular}
\caption*{\footnotesize The five pieces, mapped onto one real task.}
\end{table}

\begin{waitwhat}
The ``brain'' of the agent --- the part reasoning and writing the final
reply --- has not changed since Era 3. An agent is not a smarter model. It
is the \emph{same kind of} model, wrapped in a loop that gives it hands
(tools) and a notebook (memory). Most of this course is about building that
loop well.
\end{waitwhat}

\begin{keytake}
``Agent'' is not one capability but five, working together: reasoning,
planning, tools, memory, and autonomy. A system missing any one piece is not
yet an agent --- it is a chatbot with extra steps.
\end{keytake}

% =============================================================================
\section{The ReAct loop: the engine inside every agent}
\secsub{One repeating cycle --- reason, act, observe --- runs the five
pieces from Section 3, one step at a time.}

\vocab{ReAct} (short for \textbf{Rea}soning + \textbf{Act}ing) names the loop
almost every agent uses. The model writes a short \emph{thought} (its
reasoning). It picks one \emph{action} (a tool call). It reads the tool's
\emph{observation} (the result). Then it repeats, using each new observation
to decide the next thought, until it has enough information to answer.

\begin{intuition}
ReAct does not ask the model to solve the whole task in one breath. It asks
for one small, checkable step at a time. It must look at the real result of
that step before deciding the next one. This is the same discipline as
debugging code one line at a time, instead of guessing the whole fix at once.
\end{intuition}

\begin{everyday}
Picture solving a jigsaw puzzle with your eyes closed between moves. Thought:
``I should look for an edge piece.'' Action: pick up a piece. Observation:
open your eyes, check if it's an edge piece. If not, new thought: ``try a
different piece.'' You never plan the whole puzzle blind. You take one small
action, check the real result, and let that result shape your next move.
\end{everyday}

\housefig{figures/fig_react_loop.png}{The ReAct loop: one cycle an agent
repeats, each time using the last observation to pick the next action.}

\begin{worked}{booking the flight, four iterations of ReAct}
Request: ``Book me a window seat on tomorrow's 7am flight to Chennai, under
Rs.~6{,}000, and text me the confirmation.''
\begin{steps}
  \item \textbf{Iteration 1 --- Reason:} ``I need today's price before I can
    compare it to the Rs.~6{,}000 budget.'' \textbf{Act:} call
    \texttt{search\_flights(route=`X-Chennai', time=`7am tomorrow')}.
    \textbf{Observe:} the tool returns \texttt{\{flight: AI203, price:
    Rs.~5{,}450, seats\_left: 4\}}.
  \item \textbf{Iteration 2 --- Reason:} ``Rs.~5{,}450 is under Rs.~6{,}000, so the
    budget condition is met. Next I need a window seat.''
    \textbf{Act:} call \texttt{check\_seat\_map(flight=`AI203',
    type=`window')}. \textbf{Observe:} the tool returns \texttt{\{available:
    true, seat: `14A'\}}.
  \item \textbf{Iteration 3 --- Reason:} ``Both conditions are met: under
    budget, and a window seat is free. I can book now.'' \textbf{Act:}
    call \texttt{book\_seat(flight=`AI203', seat=`14A')}. \textbf{Observe:}
    the tool returns \texttt{\{status: `confirmed', booking\_ref:
    `PNR7X2K'\}}.
  \item \textbf{Iteration 4 --- Reason:} ``Booking is confirmed. The last
    step is to notify the user by text.'' \textbf{Act:} call
    \texttt{send\_sms(message=`Booked AI203, seat 14A, ref PNR7X2K')}.
    \textbf{Observe:} \texttt{\{status: `sent'\}}. \textbf{Final answer:}
    ``Booked! Window seat 14A on AI203, Rs.~5{,}450. Confirmation texted to
    you.''
\end{steps}
Every fact in the final answer came from a tool's real observation: the
flight number, the price, the seat, the reference. None of it came from the
model's memory. That is what ``grounded'' means.
\end{worked}

\begin{watchout}
What stops this loop from running forever on a task it cannot solve? An
\vocab{iteration limit}, set by the system's designer. For example: ``try at
most 5 times, then stop and tell the user it failed.'' Without a limit, a
confused agent can loop indefinitely. It keeps calling tools and burning
cost, with no final answer in sight. Every production agent needs this
safety rail; it is not optional.
\end{watchout}

\begin{keytake}
ReAct is reason, act, observe, repeat: one checkable step at a time. Each new
tool result shapes the very next thought. The loop runs until the agent has
enough grounded information to answer, bounded by an iteration limit so it
cannot run forever.
\end{keytake}

\noindent\textbf{Where it's used.}
ReAct-style loops power coding assistants: write code, run it, read the
error, fix it, repeat. They power research agents: search, read a result,
search again with a refined query. They power customer-support agents: look
up an order, check a policy, issue a refund. The common thread is a task
where the right next step depends on what the last action returned. It is
not a plan made in advance.

% =============================================================================
\section{System architecture, zoomed out}
\secsub{One component, the orchestrator, decides on every message whether to
think, call a tool, or check memory.}

The ReAct loop happens \emph{inside} a single agent's turn. Zoom out one
level, and a production system needs a controller for the whole
conversation. Where does incoming state live? Which model answers which kind
of message? What happens when something is slow or expensive?

\begin{everyday}
Think of a hospital reception desk. A simple question (``what are your
visiting hours?'') is answered on the spot by the receptionist --- fast,
cheap, no doctor needed. A complex case gets routed to the right specialist,
who takes longer and costs more. The receptionist is the
\vocab{orchestrator}: not the one doing the medical reasoning, but the one
deciding \emph{who} should.
\end{everyday}

\housefig{figures/fig_architecture.png}{System architecture, zoomed out: the
orchestrator decides whether a message needs the LLM's reasoning, a tool call,
or a memory lookup.}

\vocab{Conversation state} is whatever the system must remember to make
sense of the \emph{next} message. Three examples: an earlier constraint
(``under Rs.~6{,}000''), which step of a task it has reached, and the history
so far. Losing this state mid-conversation is why a chatbot
sometimes seems to ``forget'' what you just said thirty seconds ago.

\vocab{Model selection} is the trade-off between a small, fast, cheap model
and a large, slow, expensive one. A simple FAQ answer does not need the most
powerful model available. A hard, multi-step reasoning task might. Routing
simple requests to a cheap model, and escalating hard ones to an expensive
one, is a standard cost-control pattern. You will build on it directly in
later sessions on embeddings, retrieval, and evaluation.

\begin{waitwhat}
The most expensive mistake in a production agent is often not a wrong
answer. It is routing \emph{every} message, including ``what time do you
close?'', through the most powerful model available. Good orchestration
saves real money by matching effort to difficulty. A hospital does not send
every patient to a specialist, either.
\end{waitwhat}

\begin{keytake}
Above the ReAct loop sits an orchestrator. It holds conversation state
across turns, and picks which model or tool handles each message. The
choice balances latency, reliability, and cost, not just correctness.
\end{keytake}

% =============================================================================
\section*{Self-test}
\addcontentsline{toc}{section}{Self-test}

\noindent\textbf{1.} A bare LLM assistant (Era 3) can write a fluent,
confident restaurant recommendation. Why might it still be wrong?

\flipanswer{It generates from what it ``read'' during training, with no live
lookup. It can \emph{hallucinate}: state a confident claim with no tool call
to check it --- say, a restaurant that closed last year.}

\noindent\textbf{2.} Among reasoning, planning, tools, memory, and autonomy,
which one is a bare LLM assistant already good at on its own?

\flipanswer{Reasoning --- predicting good next text, including working out
what should happen next. The other four pieces (planning, tools, memory,
autonomy) turn that reasoning into an agent.}

\noindent\textbf{3.} In the ReAct loop, what directly shapes the next
Reason step, and why does that matter?

\flipanswer{The last Observe step. A new result can change the plan. If a
flight is over budget, the next thought differs from if it were under
budget. Nothing is decided blind.}

\noindent\textbf{4.} Why does a production agent need an iteration limit?

\flipanswer{Without one, a confused agent could loop forever on a task it
cannot solve, burning cost with no final answer. The limit forces a clean
stop, usually with a hand-off or an honest failure message.}

\noindent\textbf{5.} Why might a well-designed system route a simple question
to a small, cheap model instead of the most powerful one available?

\flipanswer{Cost and latency. Matching model size to task difficulty ---
the hospital-reception pattern --- avoids paying for reasoning power the
simple task never needed.}

% =============================================================================
\section*{Pocket card: Session 1 in one page}
\addcontentsline{toc}{section}{Pocket card}

\begin{keytake}
\textbf{The whole session, compressed.}
\begin{itemize}[leftmargin=1.3em,topsep=4pt,itemsep=4pt]
  \item Four eras answer ``what should the system do next?'' with growing
    capability and cost: \textbf{rules} $\to$ \textbf{retrieval} $\to$
    \textbf{generation (LLM)} $\to$ \textbf{agent loop} (generation + tools
    + memory).
  \item An \textbf{agent} needs all five of: reasoning, planning, tools,
    memory, autonomy. The model's reasoning does not change --- the loop
    around it does.
  \item The engine inside that loop is \textbf{ReAct}: reason $\to$ act (one
    tool call) $\to$ observe (the real result) $\to$ repeat, bounded by an
    \textbf{iteration limit}.
  \item Above the loop, an \textbf{orchestrator} holds conversation state
    and routes each message to the right model or tool. It trades off
    latency, reliability, and cost.
  \item \textbf{One decision rule to remember:} don't default to the most
    capable, most expensive option. Match the tool, the model, and the era
    to how hard the task actually is.
\end{itemize}
\end{keytake}

% =============================================================================
\section*{Glossary \& symbols}
\addcontentsline{toc}{section}{Glossary \& symbols}

\noindent\textbf{Key terms.}
\begin{description}[leftmargin=9em,style=nextline,font=\normalfont\bfseries\color{navy}]
  \item[chatbot] a system that replies with text only --- no tools, no
    persistent memory, no autonomous action.
  \item[agent] a chatbot wrapped in a loop with tools, memory, and bounded
    autonomy, able to act on the world, not just describe it.
  \item[reasoning] working out, from the text so far, what should logically
    happen or be said next.
  \item[planning] breaking one goal into an ordered sequence of smaller
    actions.
  \item[tool] an external capability the agent can call to act on the real
    world. Examples: an API, a search, a calculator, code execution.
  \item[memory] state carried across steps (short-term) or across sessions
    (long-term), so the agent does not re-derive context each turn.
  \item[autonomy] how much the agent may decide and do --- retry, adapt,
    stop --- without asking a human first.
  \item[hallucination] a fluent, confident claim generated by the model with
    no grounding in a real, checked source.
  \item[ReAct] the reason $\to$ act $\to$ observe $\to$ repeat loop; short
    for reasoning + acting.
  \item[iteration limit] the maximum number of ReAct cycles allowed before
    an agent must stop, solved or not.
  \item[orchestrator] the controller above the ReAct loop that holds
    conversation state and decides which model or tool handles each
    message.
  \item[model selection] choosing a small/cheap or large/expensive model per
    request, trading capability against latency and cost.
\end{description}

\medskip
\noindent\textbf{Notation at a glance.}
\begin{center}\small
\begin{tabular}{@{}ll@{}}
\toprule
\textbf{Shorthand} & \textbf{Meaning} \\
\midrule
$n$ & the iteration limit on a ReAct loop (e.g.\ $n = 5$ retries) \\
\bottomrule
\end{tabular}
\end{center}

\end{document}
